\subsection{Atomic Memory-Order Selectors}
\Needspace{1.25in}
\BedrockTableCaption{Atomic Memory-Order Selectors}
\begingroup\footnotesize
\setlength{\aboverulesep}{0pt}
\setlength{\belowrulesep}{0pt}
\setlength{\extrarowheight}{1.2pt}
\begin{center}
\begin{tabularx}{\linewidth}{@{}>{\raggedright\arraybackslash}p{0.95in}>{\raggedright\arraybackslash}p{0.45in}>{\raggedright\arraybackslash}X@{}}
\toprule
\rowcolor{BedrockHeaderFill}
\textbf{Selector} & \textbf{Code} & \textbf{Architectural Ordering Effect}\\
\midrule
\texttt{RELAXED} & \texttt{0} & Atomicity only. No ordering is imposed on unrelated (:term:base.terms.memory_operation|plural:).\\
\texttt{ACQUIRE} & \texttt{1} & Later (:term:base.terms.memory_operation|plural:) by the same (:term:base.terms.logical_processor:) are ordered after the atomic operation.\\
\texttt{RELEASE} & \texttt{2} & Earlier (:term:base.terms.memory_operation|plural:) by the same (:term:base.terms.logical_processor:) are ordered before the atomic operation.\\
\texttt{ACQREL} & \texttt{3} & Applies both acquire and release ordering.\\
\texttt{SEQCST} & \texttt{4} & Applies acquire-release ordering and participates in the single global sequentially consistent order.\\
\bottomrule
\end{tabularx}
\end{center}
\endgroup


(:ref:base.instructions.CMPXCHG:), (:ref:base.instructions.FETCHADD:), (:ref:base.instructions.FETCHSUB:), (:ref:base.instructions.FETCHAND:), (:ref:base.instructions.FETCHOR:), and
(:ref:base.instructions.FETCHXOR:) are the atomic read-modify-write instructions. Each performs its memory read, successful memory
write, and architectural register result as one indivisible operation in the coherence order of the addressed
physical location.

A successful release or stronger atomic write heads a release sequence. The sequence continues through immediately
following successful atomic read-modify-write operations to the same location in coherence order and ends before any
other write. An acquire or stronger atomic operation that takes its value from any member of the sequence observes
the effects ordered before its head before effects ordered after the acquiring operation.

One encoded selector governs the complete (:ref:base.instructions.CMPXCHG:) operation on both success and failure. A failed
(:ref:base.instructions.CMPXCHG:) contributes a memory read event. Its acquire component orders later (:term:base.terms.memory_operation|plural:) after that read, and
its release component orders earlier (:term:base.terms.memory_operation|plural:) before that read. On success, the memory write event carries the
selected release effect to observers of that write; on failure, the read is the operation\textquotesingle{}s (:term:base.terms.memory_event:). A failed
(:ref:base.instructions.CMPXCHG:) contributes no write and therefore creates no release sequence. A failed
\texttt{CMPXCHG.SEQCST} additionally participates in the global sequentially consistent order as an SC load.
